%% Source: MANUSCRIPT.md "Materials and Methods > Reaction similarity"
%%         (drafts/methods_reaction_similarity.md).
\subsection{Reaction similarity}\label{sec:methods-similarity}

Reaction similarity --- a per-pair distance or similarity score across all
mass-balanced ModelSEED reactions --- is a new capability in the 2026 release.
It supports two use cases: as a diagnostic for the reconciliation ontology
described in the 2020 paper (reactions with high similarity that are not already
ontology-linked are candidates for review), and as an input feature for
downstream applications including annotation transfer, gap-filling, and
pathway-inference.

\paragraph{Pipeline.} Every ModelSEED reaction is represented by the
concatenation of its balanced reactant and product structures. Structures are
passed through an external reaction-embedding foundation model to produce a
fixed-dimensional vector; the similarity between two reactions is the cosine
similarity between their embeddings. GPU-accelerated batching allows the full
pairwise similarity matrix to be regenerated in under approximately ten minutes
on a single GPU, so the matrix is treated as a derived artifact and refreshed
whenever compound curation materially changes the input structures.

\paragraph{External reaction-embedding foundation-model selection.} Several
external reaction-embedding models are available at the time of this release.
Rather than adopting one by default, we evaluate the candidate models on their
discerning power over ModelSEED reactions and defend the chosen model on that
basis. The evaluation compares candidate models on: (i) intra-EC clustering ---
reactions sharing the same EC number should embed close together; (ii)
intra-pathway clustering --- reactions within a KEGG or MetaCyc pathway should
embed close together; (iii) retrieval quality on a held-out set of known
reaction-pair relationships (e.g. adjacent reactions in a pathway). Candidate
models compared: \TBD[list]. Chosen model: \TBD{} on the basis of
\TBD[metric result].

\paragraph{Publication of the matrix.} The similarity matrix is exported as a
compressed sparse artifact under \file{Biochemistry/Reaction\_Similarity/}. Users
can regenerate it locally by rerunning the pipeline script against a specific
ModelSEED release tag.
